\documentclass[book.tex]{subfiles}

\begin{document}
\chapter{Fermions and Bosons}

Particles come in two fundamental types, fermions and bosons. These two particles are differentiated by whether they obey the Pauli exclusion principle, which states that two particles with half-integer spins cannot occupy the same quantum state.

\section{Statistics of Particles}

As we briefly mentioned, fermions and bosons obey different statistical properties.

\subsection{Bose-Einstein Statistics}

Bosons are indistinguishable from one-another and can occupy the same energy state. The probability density collection for a set of bosons is given by

\begin{align}
f(E_i) = \frac{g_i}{e^{(E_i - \mu)/k_BT} - 1},
\end{align}
where $g_i$ is the degeneracy (the number of quantum states that share the same energy level), $k_B$ is Boltzmann's constant, $E_i$ is the energy of the $i$-th possible state, $T$ is the absolute temperature, and $\mu$ is the chemical potential. This equation can be derived from the form of the grand canonical ensemble that we encountered earlier in the book.

\subsection{Fermi-Dirac Statistics}

The Fermi-Dirac statistics describe the distributions ion over energies for an ensemble of particles that obey the Pauli exclusion principle and which is at thermal equilibrium. As a technical condition, these particles have half integer spin and are called fermions. The energy distribution is given by the following probability density function.

\begin{align}
    f(E_i) = \frac{1}{1 + e^{(E_i - \mu)/k_BT}},
\end{align}

This formula can also be derived from the grand canonical ensemble. Electrons are examples of fermions.

 %(TODO: Add a description of the chemical potential to the statistical mechanics chapter.)


\end{document}